\section{A classification in degree five}
\label{sec:quintic}

\textbf{Theorem.} Across all levels and nebentypus, exactly twenty-six
quintic fields occur as coefficient fields of weight-two CM newforms.
One is the maximal real subfield of $\mathbb Q(\zeta_{11})$. Each of
the other twenty-five comes from a different imaginary quadratic field
of class number five. This result is unconditional. Martin singled out
degree five in Remark 1.2 as a possible extension of his explicit
classification. \href{https://arxiv.org/pdf/2509.24119v2}{Martin, Remark 1.2}.

\paragraph{The input from degree five.}
As in the cubic argument, the odd coefficient degree implies trivial
nebentypus. Put $L=EK_f$, so $[L:E]=5$ and $\exp\mathrm{Cl}(E)\mid5$.
For $h_E>1$, $E$ has only the units $\pm1$.
Let $r$ be the order of the finite character $\eta$.
As $\varphi(r)\mid10$ and $r$ is even, $r$ is one of $2,4,6,22$.
Orders four and six would force $E$ to be the quadratic cyclotomic
field inside $L$, since $[L:E]$ is odd. Order twenty-two would
force $L=\mathbb Q(\zeta_{11})$ and hence $E=\mathbb Q(\sqrt{-11})$.
All three exceptional quadratic fields have class number one.
Thus $r=2$.

For each class $[\mathfrak a]$, use a representative prime to the
conductor, and write $\mathfrak a^5=(\theta)$. The classes $[\theta]$
form a subgroup of $E^\times/(E^\times)^5$ isomorphic to
$\mathrm{Cl}(E)$ by the ideal argument of the preceding section.
Each $\theta$ has a fifth root in $L$ since $\psi(\mathfrak a)^5=\pm\theta$.
These classes stay independent over $F=E(\zeta_5)$, as $[F:E]$ divides
four. Kummer theory gives $h_E\mid[FL:F]=5$. The only cases to
consider are therefore $h_E=1$ and $h_E=5$.

\paragraph{Class number one.}
Principal ideals supply all values, so $L=E(\mu_r)$.
For this field to have degree ten over $\mathbb Q$ the preceding
cyclotomic argument leaves only $E=\mathbb Q(\sqrt{-11})$ and
$r=22$. To realize it, reduce modulo $11\mathcal O_E$. The quotient
is the ring of dual numbers over $\mathbb F_{11}$, with unit group
$\mathbb F_{11}^\times\times(1+\epsilon\mathbb F_{11})$,
$\epsilon^2=0$. Choose the Legendre character on the first factor
and a character of order eleven on the second. The resulting
$\eta$ sends $-1$ to $-1$ and restricts on rational units to
$\chi_{-11}$. Hence $\psi((\alpha))=\alpha\eta(\alpha)$ defines a
type-one Hecke character, primitive because the order-eleven part
is nontrivial. Its newform has trivial nebentypus and the asserted
real quintic coefficient field.

\paragraph{Class number five.}
Genus theory and the odd class number give a single prime discriminant,
necessarily $D=-q$ with $q\equiv3\pmod4$ an odd prime; the purely
dyadic cases have class number one. Let $\mathfrak t$ be an integral
ideal representing a generator, prime to $q$, with norm $n$.
Choose $\theta$ with $\mathfrak t^5=(\theta)$.
There is exactly one possible coefficient field up to isomorphism.
Indeed for any conductor take a representative $\alpha\mathfrak t$
coprime to it. Its character value, divided by $\alpha$, is a fifth
root of $\pm\theta$. The principal-ideal values generate $E$, and this
one value generates the rest of $L$ over $E$. The sign can be absorbed
in the choice of fifth root. The polynomial $X^5-\theta$ is
irreducible over $E$: if a root generated an extension of degree
$1\leq j<5$, its norm would show $\theta^j$ is a fifth power, forcing
$\theta$ itself to be a fifth power. This contradicts the order of
$\mathfrak t$. Thus all choices yield isomorphic $L$ over $E$.
As $L$ is a totally imaginary quadratic extension of its totally real
coefficient field, its maximal totally real subfield recovers $K_f$.

For existence, choose conductor the prime above $q$ and $\eta$ the
Legendre character on its residue field. It sends $-1$ to $-1$ and its
rational restriction is $\chi_D$. The principal-ideal rule
$\psi((\alpha))=\eta(\alpha)\alpha$ extends to all ideals prime to the
conductor by assigning to $\mathfrak t$ any fifth root of
$\eta(\theta)\theta$: every ideal is a principal ideal times a power
of $\mathfrak t$, and its only generator relation modulo principal
ideals is the fifth-power relation used here. This character is primitive
at $q$. Its newform has weight two, trivial nebentypus and coefficient
degree five. If $z=\psi(\mathfrak t)$ and $\Theta=z^5$, then
$\overline z=n/z$ because $|\theta|^2=n^5$. In particular the field
is generated by $x=z+n/z$, a root of
\[
X^5-5nX^3+5n^2X-\operatorname{Tr}_{E/\mathbb Q}(\Theta).
\]
Indeed $x$ is real and $[E(z):E]=5$, while $z$ satisfies
$Z^2-xZ+n=0$, so $x\notin\mathbb Q$.

\paragraph{Why different CM fields give different quintic fields.}
Set $B=E(\zeta_5)$ and $M=B(z)$. Since $[B:E]$ divides four,
$\Theta$ remains a non-fifth-power in $B$. Thus $M/B$ is cyclic of
degree five. Conjugation of $\Theta$ takes it to $n^5/\Theta$, so $M$
is normal over $\mathbb Q$. It is unramified outside $5q$: at all
other places $B$ is unramified and the valuation of $\Theta$ is a
multiple of five. Locally divide by a fifth power of a uniformizer
and adjoin a fifth root of a unit. Its reduction is separable at that
residue characteristic, giving an unramified extension.

At $q$, an inertia element restricting nontrivially to $E$ fixes
$\zeta_5$ and sends $z$ to $\zeta_5^j n/z$ for some $j$. The five
conjugates of $x$ are $x_i=\zeta_5^i z+\zeta_5^{-i}n/z$.
They are distinct, since an equality for two different $i$ would make
$z^2/n$ a root of unity, contrary to the different ideals
$\mathfrak t^5$ and $\overline{\mathfrak t}^{5}$.
Indeed equality of these ideals would imply $[\mathfrak t]=[\mathfrak
t]^{-1}$. The inertia element acts on the five conjugates by
$i\mapsto-i-j$ and therefore acts nontrivially. This forces
ramification in the quintic field: if it were unramified at $q$, the
compositum of all its local conjugates would be unramified there as
well. Thus $q$ is the unique ramified rational prime other than five
and recovers the CM field. The cyclotomic example ramifies only at
eleven, and $\mathbb Q(\sqrt{-11})$ has class number one, so it is
distinct from every class-number-five example.

The unconditional class-number-five classification contains twenty-five
fields, with largest absolute discriminant $2683$. This proves the
count. \href{https://magma.maths.usyd.edu.au/~watkins/papers/smallcn.pdf}{Watkins, \emph{Class numbers of imaginary quadratic fields}, final table}.

The canonical form of level $q^2$ and coefficient degree $h(-q)$ is
part of earlier work, including Gonz\'alez--Lario, Proposition~5.3 and Gross, Section~8.
Yang's class-number divisibility gives the reduction to class number one or
five directly. The additional assertions here are exhaustiveness over all
levels, distinction by the ramified prime, and the explicit polynomial list.
\href{https://arxiv.org/pdf/0810.4224}{Gonz\'alez and Lario, Proposition~5.3};
\href{https://people.math.harvard.edu/~gross/preprints/hecke.pdf}{Gross, Section~8}.
